\chapter{کلیات پروژه}
\label{chap:introduction}

\section{مقدمه}
گسترش منابع انرژی تجدیدپذیر و ذخیره‌سازهای الکتریکی، ساختار سنتی شبکه قدرت را به سمت سامانه‌های توزیع‌شده و هوشمند سوق داده است. ریزشبکه مجموعه‌ای از منابع تولید پراکنده، بارهای محلی و تجهیزات ذخیره‌سازی است که می‌تواند در حالت متصل به شبکه اصلی یا به‌صورت جزیره‌ای کار کند. بهره‌برداری مناسب از این اجزا به یک سامانه مدیریت انرژی نیاز دارد تا در هر بازه زمانی، میزان تولید، مصرف، شارژ یا دشارژ باتری و تبادل توان با شبکه اصلی را تعیین کند.

مدیریت انرژی ریزشبکه تنها یک مسئله پیش‌بینی نیست؛ تصمیم‌های حاصل باید محدودیت‌های فیزیکی مانند توازن توان، ظرفیت باتری، حدود توان مبدل‌ها و دامنه مجاز وضعیت شارژ را رعایت کنند. از طرف دیگر، عدم قطعیت تولید خورشیدی و تغییرات بار سبب می‌شود که مدل‌سازی دقیق سامانه دشوار باشد. روش‌های بهینه‌سازی کلاسیک می‌توانند این قیود را به‌طور صریح اعمال کنند، ولی حل مکرر آن‌ها در مسائل بزرگ یا برخط ممکن است پرهزینه باشد. روش‌های یادگیری ماشین سرعت استنتاج مناسبی دارند، اما اگر صرفاً از داده استفاده کنند، لزوماً سازگاری فیزیکی خروجی را تضمین نمی‌کنند.

شبکه‌های عصبی فیزیک‌آگاه یا \lr{PINN} روشی برای ترکیب داده و دانش پیشین فیزیکی فراهم می‌کنند. در این رویکرد، معادلات حاکم بر سامانه بخشی از تابع هزینه آموزش هستند؛ بنابراین شبکه عصبی افزون بر برازش داده، برای رعایت قوانین فیزیکی نیز آموزش می‌بیند. چارچوب عمومی این روش توسط رئیسی، پردیکاریس و کارنیاداکیس برای حل و شناسایی معادلات دیفرانسیل غیرخطی ارائه شده است \pcite{raissi2019pinns}.

در سوی دیگر، پژوهش‌های مدیریت انرژی ریزشبکه نشان می‌دهند که مدل بهره‌برداری باید هم‌زمان تبادل انرژی، سود یا هزینه مشارکت‌کنندگان و قیود تجهیزاتی را در نظر گیرد. برای نمونه، اولادجو و فولی یک مدل ریزشبکه متصل به شبکه را با تولید خورشیدی، باتری و چند مشارکت‌کننده بررسی کرده و برای تخصیص منصفانه سود از نظریه بازی و برای حل مسئله از الگوریتم بهینه‌سازی مبتنی بر آموزش و یادگیری استفاده کرده‌اند \pcite{oladejo2019microgrid}. پروژه حاضر از ترکیب ایده فیزیک‌آگاه با ساختار فیزیکی و اقتصادی مدیریت انرژی ریزشبکه الهام می‌گیرد.

\section{بیان مسئله}
یک مدل یادگیری برای مدیریت انرژی ریزشبکه، نگاشتی میان اطلاعات ورودی و تصمیم‌های بهره‌برداری ایجاد می‌کند. ورودی‌ها می‌توانند شامل زمان، تقاضای بار، توان خورشیدی، قیمت برق و وضعیت شارژ باتری باشند. خروجی‌ها نیز می‌توانند توان دریافتی از شبکه، توان شارژ و دشارژ باتری یا برنامه تبادل انرژی را نشان دهند. اگر این نگاشت تنها با کمینه‌سازی خطای داده آموزش داده شود، ممکن است خروجی مدل از نظر عددی به داده نزدیک باشد، اما معادله توازن توان یا حدود ذخیره‌ساز را نقض کند.

مسئله اصلی این پروژه آن است که چگونه می‌توان دانش فیزیکی ریزشبکه را در فرایند آموزش شبکه عصبی وارد کرد تا مدل، علاوه بر دقت داده‌ای، سازگاری فیزیکی بیشتری داشته باشد. ایده محوری آن است که باقیمانده معادلات توازن و پویایی ذخیره‌ساز در نقاط آموزش محاسبه و به تابع هزینه افزوده شود. به این ترتیب، فضای پاسخ‌های قابل قبول شبکه عصبی محدودتر می‌شود و انتظار می‌رود مدل با داده کمتر نیز رفتار معقول‌تری داشته باشد.

\section{اهمیت و ضرورت}
استفاده از مدل‌های فیزیک‌آگاه در ریزشبکه از چند جهت اهمیت دارد:
\begin{itemize}
  \item داده‌های واقعی سامانه قدرت ممکن است محدود، پرهزینه یا همراه با نویز باشند؛
  \item قوانین اصلی مانند بقای انرژی از پیش معلوم‌اند و نادیده‌گرفتن آن‌ها موجب اتلاف اطلاعات می‌شود؛
  \item در کاربردهای کنترلی، نقض قیود فیزیکی می‌تواند تصمیم تولیدشده را غیرقابل اجرا کند؛
  \item یک مدل آموزش‌دیده می‌تواند پس از آموزش، تصمیم یا تخمین را با سرعت بیشتری تولید کند؛
  \item ترکیب مدل‌محور و داده‌محور امکان بررسی تعادل میان دقت، سرعت و قابلیت اعتماد را فراهم می‌سازد.
\end{itemize}

\section{اهداف پروژه}
هدف اصلی، تدوین و پیاده‌سازی یک چارچوب اولیه برای کاربرد شبکه‌های عصبی فیزیک‌آگاه در مدیریت انرژی ریزشبکه است. اهداف فرعی عبارت‌اند از:
\begin{enumerate}
  \item مطالعه مبانی شبکه‌های عصبی فیزیک‌آگاه و مدیریت انرژی ریزشبکه؛
  \item استخراج معادلات و قیود مناسب برای یک ریزشبکه متصل به شبکه؛
  \item تعریف تابع هزینه ترکیبی شامل مؤلفه داده و مؤلفه فیزیکی؛
  \item طراحی ساختار قابل پیاده‌سازی با پایتون و پای‌تورچ؛
  \item تعریف یک روش ارزیابی برای مقایسه مدل فیزیک‌آگاه و مدل صرفاً داده‌محور؛
  \item سنجش خطای پیش‌بینی، باقیمانده فیزیکی، نقض قیود و هزینه بهره‌برداری پس از اجرای آزمایش‌ها.
\end{enumerate}

\section{پرسش‌های پروژه}
این پروژه بر پرسش‌های زیر متمرکز است:
\begin{enumerate}
  \item چگونه می‌توان معادلات توازن توان و وضعیت شارژ باتری را در تابع هزینه شبکه عصبی وارد کرد؟
  \item وزن مؤلفه‌های داده‌ای و فیزیکی تابع هزینه چه اثری بر آموزش مدل دارد؟
  \item آیا مدل فیزیک‌آگاه نسبت به مدل صرفاً داده‌محور، نقض قیود کمتری ایجاد می‌کند؟
  \item چه معیارهایی برای ارزیابی هم‌زمان دقت و سازگاری فیزیکی مناسب‌اند؟
\end{enumerate}

\section{دامنه و محدودیت‌ها}
مدل اولیه این گزارش یک ریزشبکه متصل به شبکه شامل تولید خورشیدی، باتری، بار محلی و تبادل توان با شبکه اصلی را در نظر می‌گیرد. تمرکز اصلی بر لایه مدیریت انرژی است و جزئیات الکترومغناطیسی مبدل‌ها یا پایداری گذرای شبکه در دامنه نسخه اولیه قرار نمی‌گیرد. همچنین نتایج عددی باید بر اساس داده و اجرای واقعی کد گزارش شوند و هیچ نتیجه‌ای پیش از انجام آزمایش فرض نمی‌شود.

\section{ساختار گزارش}
در فصل~\ref{chap:background} مبانی ریزشبکه، مدیریت انرژی و شبکه‌های عصبی فیزیک‌آگاه مرور می‌شود. فصل~\ref{chap:method} چارچوب پیشنهادی و تابع هزینه را توضیح می‌دهد. فصل~\ref{chap:implementation} به ساختار پیاده‌سازی و روش ارزیابی اختصاص دارد. سرانجام، فصل~\ref{chap:conclusion} جمع‌بندی و مسیرهای ادامه کار را ارائه می‌کند.
